\chapter{Assembly Index and the Depth of the Fractal Learner}
\label{sec:assembly-depth}
% ===================================================================

Two constructions in this book measure the same thing and have not yet
been introduced to each other.

Chapter~\ref{sec:assembly} identified the assembly index of a term with
the minimum path length in its open synchronization tree: the fewest
environment-assisted joining steps that suffice to build it from
elementary material.
Chapter~\ref{ch:eng} built a tower --- the outside of one level is the
inside of the next --- and observed that an individual is a fixed point
of a partition whose multiple solutions are the levels of that tower.
The first is a number attached to an object by a chemist.
The second is a structural fact about a mind assembled out of learners.

This chapter argues that the second is bounded by the first, and draws
the consequence: assembly theory, taken seriously, gives an upper bound
on what an ecology-as-mind can know.

\section{Every level of the tower costs a joining step}

Recall the generating step of Chapter~\ref{ch:eng}.
A \emph{medium} is a channel over which the parts of a region
communicate, and $@C$ --- the quotation of a medium --- sits one level up
the reflection tower from the names it carries, because $C$'s free names
are the endpoints it joins.
Proposition~\ref{eng:prop:tower} says that the media of one level are the
channels of the next, so ascending the tower means constructing, at each
level, a carrier for the names of the level below.

That construction is not free, and its cost is exactly the kind of cost
the assembly index counts.

\begin{proposition}[Tower height is bounded by assembly index]
\label{prop:tower-ai}
Let $E$ be an ecology assembled from elementary terms, and let $h(E)$ be
the height of its medium tower in the sense of
Proposition~\ref{eng:prop:tower}.
Then
\[
  h(E) \;\leq\; \AI(E),
\]
where $\AI$ is the assembly index of Proposition~\ref{prop:assembly-index}.
\end{proposition}

\begin{proof}[Proof sketch]
Ascending one level of the tower requires exhibiting a medium $C$ whose
endpoints are names of the level below.
By the no-self-code theorem a name codes a strictly smaller term, so $C$
is not among the terms it carries and must be constructed.
Its construction is a transition in the open synchronization tree: it
joins previously present material at a minimal context, which is what an
edge of $\mathcal{ST}_o$ is.
Hence each level contributes at least one edge to any path from the
elementary terms to $E$, and in particular to a minimum-length such path.
Summing over levels gives $h(E) \leq \AI(E)$.
\end{proof}

\begin{remark}[The bound is loose, and interestingly so]
\label{rmk:ai-loose}
The inequality is not tight, and the slack is where the biology is.
An organism may spend a great many joining steps on structure that
contributes no new level --- redundancy, repair machinery, sheer bulk ---
and Proposition~\ref{eng:prop:radius} says it must spend some, because a
medium beyond the profitable radius carries no traffic and a level with
no traffic is not a level.
So $\AI$ counts everything built and $h$ counts only the building that
bought a new inside.
The ratio $h(E)/\AI(E)$ is therefore a measure of architectural
efficiency, and one would expect it to be small for a crystal, larger for
a bacterium, and larger still for a nervous system.
We do not know how to compute it for any real organism, and we note that
the quantity is at least in principle measurable, which is more than can
be said for most proposed complexity measures.
\end{remark}

\section{Depth bounds what can be told apart}

The bound becomes interesting when combined with what the learner's
hypothesis language can express.

Chapter~\ref{ch:sci} established that a mortal scientist's hypotheses are
formulae of a logic of \emph{namespaces} rather than of individuals,
because a specimen affords exactly one measurement and the measurement
ends it.
Chapter~\ref{sec:lts-hml} gave that logic its modal operators, indexed by
minimal contexts.
The two facts together mean that a hypothesis is a modal formula whose
nesting alternates over namespace boundaries.

\begin{proposition}[Affordable nesting is bounded by tower height]
\label{prop:nesting-bound}
A learner in an ecology $E$ can form and evaluate hypotheses of namespace
nesting depth at most $h(E)$.
\end{proposition}

\begin{proof}[Proof sketch]
A formula of nesting depth $d$ quantifies over $d$ distinguishable
namespace levels.
Evaluating it requires access to a percept that has traversed those
levels, which by Definition~\ref{eng:integration-depth} is a percept of
integration depth $d$.
By Proposition~\ref{eng:prop:fragile} such a percept survives with
probability $p^d$ and costs $d$ metered rendezvous.
If $d$ exceeds $h(E)$ there are no further levels for it to traverse, and
the formula distinguishes nothing the shallower formulae did not.
\end{proof}

Now recall the standard fact about modal logic that makes this bite:
formulae of Hennessy--Milner nesting depth $n$ characterise exactly
$n$-step bisimulation.
Two processes agreeing on all formulae of depth $n$ are indistinguishable
by any experiment of $n$ rounds, however cleverly designed.

\begin{theorem}[Assembly index bounds discriminating power]
\label{thm:ai-bound}
Let $E$ be an ecology-as-mind with assembly index $\AI(E)$.
Then $E$ cannot distinguish two environments that are
$\AI(E)$-step bisimilar.
Equivalently, the ontology available to $E$ is at best the quotient of
its world by $\AI(E)$-step bisimulation.
\end{theorem}

\begin{proof}
Compose Proposition~\ref{prop:tower-ai}, Proposition~\ref{prop:nesting-bound},
and the Hennessy--Milner characterisation.
\end{proof}

\begin{remark}[What ``computational power'' can and cannot mean here]
\label{rmk:power-caveat}
Theorem~\ref{thm:ai-bound} is not a statement about Turing degree, and it
would be a serious misreading to take it as one.
The rho calculus is Turing complete, and so is every ecology built in it,
at every assembly index above the trivial: a single sufficiently clever
term computes everything computable, given unbounded time and unbounded
tokens.
Assembly index does not bound what $E$ could compute in principle.

What it bounds is what $E$ can \emph{tell apart}, and that is the notion
of power the rest of this book has been using.
Chapter~\ref{ch:sci} showed that the equivalence a budgeted learner can
afford is coarser than the one that is true, and that the gap is
economics rather than ignorance.
Theorem~\ref{thm:ai-bound} says how much coarser, in terms of a quantity
a chemist can in principle measure.
The bound is on resolution, not on reach.
\end{remark}

\section{Copy number is the other axis}

Assembly theory carries two parameters, and so far only one has been
used.
The second is copy number, given a rigorous definition in
Definition~\ref{def:copy-number} as the count of bisimilar processes
sitting at channels within a namespace region.

Copy number does not extend the depth bound; it does something
orthogonal.
Where assembly index bounds how many levels the tower has, copy number
bounds how many independent samples the learner can take at any level ---
which is to say, how much statistical confidence it can buy.
Chapter~\ref{ch:sci} required this without naming it: hypotheses are about
kinds because a specimen affords one measurement, and confidence about a
kind requires many specimens of it.

\begin{remark}[Depth and width price differently]
\label{rmk:depth-width}
The two parameters have different economics, and the difference is
familiar from Chapter~\ref{ch:eng}.
Depth is exponentially fragile: reliability $p^d$, so each additional
level costs geometrically more in retransmission and repair.
Width is linear: twice the copies, twice the samples, twice the tokens.
An ecology under resource pressure will therefore trade depth away first,
and one predicts --- as an empirical matter, not a theorem --- that
lineages entering a period of scarcity should lose architectural levels
before they lose population.
Assembly theory's joint claim, that high assembly index \emph{together
with} high copy number is the biosignature, reads in this framework as
the observation that only a system with a metabolism can afford to buy
both at once.
\end{remark}

\section{Apertures, and the connection to choice}

Chapter~\ref{ch:ent} left an obligation, and this is the place to
discharge it.

The argument there was that whatever exceeds Turing computation in a
physical system must be spent at the apertures --- the cuts where
contention is observable and a scheduler must pick.
It was then asserted that the number of apertures grows with the
structural elaboration of a population rather than with its raw size, and
that the thing which measures elaboration is the assembly index.

Proposition~\ref{prop:tower-ai} is what makes that precise.
Apertures live at boundaries between levels, because a cut internal to a
level with a single owner is resolved by that owner and carries no
observable choice.
So the count of aperture classes is bounded by the number of level
boundaries, which is $h(E) - 1$, which is bounded by $\AI(E)$.

\begin{corollary}[Assembly index bounds available choice]
\label{cor:ai-choice}
The number of distinguishable aperture classes in an ecology $E$ is at
most $\AI(E) - 1$.
Consequently, if Conjecture~\ref{conj:dissipation-choice} holds, the
choice strength an ecology can exhibit is bounded by a function of its
assembly index.
\end{corollary}

This is the sharpest form of the framework's claim about mind, and it is
worth stating plainly.
A thing that has been assembled by $n$ joining steps can have at most $n$
levels of inside, can afford hypotheses of at most $n$ nested
namespaces, can resolve its world at best to $n$-step bisimulation, and
can host at most $n-1$ kinds of genuine indeterminacy.
Mind is bounded by manufacture.

\section{Open gaps}

Three, and they are not small.

The first is that Proposition~\ref{prop:tower-ai} identifies a level of
the medium tower with at least one joining step, but does not show that
the joining steps for distinct levels are distinct edges of a single
minimum-length path.
If a single construction could serve two levels the count would collapse.
We believe it cannot, on the strength of the no-self-code theorem, but
have not written the argument.

The second is that Proposition~\ref{prop:nesting-bound} treats namespace
nesting depth and modal nesting depth as the same number.
They are not obviously the same number.
The identification is exactly the sort of thing the OSLF construction of
Chapter~\ref{ch:hypercube} should settle, and until the manufacture of
the logic is written down rather than described, the identification is a
plausible assumption rather than a fact.

The third is that assembly index is defined for an object and an ecology
is not obviously an object.
Chapter~\ref{ch:eng} argued that an individual is a fixed point of a
partition, and that argument gives us the right to speak of $\AI(E)$ only
once a solution to the fixed-point condition has been chosen.
Different solutions give different indices, and we have said nothing
about how they compare.
This is the same circularity that made individuation a fixed point rather
than a construction, arriving now in a place where it costs a number.

% ===================================================================
